\documentclass[12pt]{article}

\usepackage{float}
\usepackage{graphicx} % Required for inserting images
\usepackage{fontspec}
\usepackage{apacite}
\setmainfont{Times New Roman}
\usepackage[letterpaper,top=2cm,bottom=2cm,left=2cm,right=2cm,marginparwidth=1.5cm]{geometry}

\title{Weekly assignment 4: Legal aspects}
\author{Truong Khang Tobias Nguyen}
\date{}

\begin{document}

\maketitle

\section{What I learned}
Coming from a technical background, I did not have any prior exposure to the legislative aspects of cybersecurity and thus, I (obviously) expected to learn more about it and considering that we live in Europe, I expected the main focus of this lecture to be the GDPR. 
\newline
That being said, what I learned from the lecture were the fundamentals of the GDPR, it's use, applications as well as the basics to concepts such like privacy itself. What I particularly enjoyed from the lecture were the discussions for the StudentPulse platform that we had after every logical section of the lecture. That way, Jens made sure that we were not only just sitting out the lecture but also think about the concepts he explained and apply them to a practical use case.
\newline
What I also thought was really good was that although Jens thoroughly explained the most important concepts of the GDPR but did not go into too much depth in a way that would have required us to look up or memorize paragraphs of the GDPR.
\newline
Additionally, as the GDPR introduces terms or principles such as purpose limitation, special categories, data minimization, etc. whose definitions may or may not be instantly known, we were also taught was what those terms mean and why they are important.
\newline
To conclude this section, what I learned are the fundamentals of the GDPR, it's underlying concepts and how it is applied in organizations. The lecture also taught me to reason about the requirements that need to be met so that organizations can also comply to the GDPR and the elements they have to consider as well as the possible consequences when these requirements are not met or when organizations do not comply with them.

\section{What I missed}
Although the lecture explained the most relevant aspects of the GDPR really well, there could have been additional concepts taught or explained to provide us with more context and subsequently, reinforce our understanding. 
\begin{itemize}
\item \textbf{Historical context:} Although the introduction to the GDPR was very extensive, what I would have liked to be included would be the historical background of the GDPR, e.g. what were the main drivers of it's implementation, why it was implemented, or if there was any kind of incident that led the EU to formulate it. With this context, I believe that it would be easier for us to understand why the GDPR exists and why it is so important.
\item \textbf{Counterparts:} What I would have been really interested in, would have been a slight introduction to counterparts of the GDPR. As the GDPR only exists in Europe, I am interested in their counterparts in different regions in the world such as the US or China, as well as the differences between them and their implications on organizations that operate internationally and thus, have to follow multiple laws.
\item \textbf{Critical perspectives:} Although I believe that the GDPR is important and that the information of individuals need to protected, I would also like to know different views on the GDPR as there are many critical points to it such as it's high compliance and administrative costs, inconsistent enforcement between the member states and barriers to innovation. I would have liked to know how these criticism came to place, whether the GDPR needs to be adapted according to them and the implications of a softer GDPR would be.
\item \textbf{Organizational implications:} Another point that I would have liked to hear about is how the GDPR itself affects organizations in how they operate as well as their cybersecurity culture. Despite the GDPR being a law that organizations need to comply to, it is interesting to know the differences between organizations that have a high acceptance to the GDPR and organizations that have a low acceptance of it and subsequently only follow it to not be punished. One aspect I am interested about is whether a high acceptance of the GDPR also leads to a more mature security landscape inside an organization.
\end{itemize}
Although these points were missing to me in the lecture, I also understand that the lecture lasted for four hours and with limited time, some topics either have to be slimmed down or be cut entirely. If these points were introduced in the lecture, it would have led to the other aspects of the lectures not be sufficiently taught.

\section{What I would've done differently}
Considering that I have enjoyed the lecture, there isn't much that I would have done differently.
\newline
As said, although I would have liked for the missed topics to have also been taught, that would have led to the other content of the lecture to not be adequately taught, so in that sense I accept that they were not taught.
\newline
Other than that, I believe that we could have been given more time on the discussions on the StudentPulse platform as we had around 2-3 minutes to discuss them and present our findings. Although not a bad thing by itself, it was difficult to get meaningful results when the discussions were about paragraphs or attributes that we had to map to the discussion. As the GDPR is big, it takes a while to find the right resources in the first place before one can put any thought to them. That being said, I would have given more time on the discussions or at least given more time on those that required research.
\newline
Also, I would have moved the Kahoot quiz to the end of the lecture. Although I understand that the goal of the quiz was to get a feeling of our existing knowledge of the topic, I am unfamiliar with the topic and subsequently had difficulties with the quiz. I think that moving the quiz to the end of lecture would have the effect of binding the knowledge that was taught during the course of the lecture.
\newline
Finally, I would have designed the lecture in way that would have left around 10 minutes at the end for students to ask questions. Because all the topics were condensed into a 4-hour window, it left no time for discussions afterwards which forced me to ask Jens after the lecture has actually ended.

\end{document} 